\documentclass[spanish,letterpaper,oneside]{article}

\def\documenttitle {Álgebra Lineal}
\def\documentsubtitle {Portafolio de evidencias}
\def\documentsubject {Ingeniero Agrónomo}
\def\documentauthor {Integrantes por registrar}
\def\coursename {Álgebra Lineal}
\def\coursecode {Clave por registrar}
\def\universityname {Universidad Autónoma de Nuevo León}
\def\universityfaculty {Facultad de Agronomía}
\def\universitydepartment {Ingeniero Agrónomo}
\def\universitydepartmentimage {UANL/assets-uanl/logo-uanl.png}
\def\universitydepartmentimagecfg {width=3.2cm}
\def\universitylocation {General Escobedo, Nuevo León, México}
\def\authortable {
  \begin{tabular}{ll}
    \textbf{Integrantes:} & Por registrar \\
    Matrículas: & Por registrar \\
    Grupo: & Por registrar \\
    Programa: & Ingeniero Agrónomo \\
    Unidad de aprendizaje: & \coursename \\
    Docente: & Por registrar \\
    & \\
    \multicolumn{2}{l}{Fecha de realización: \today} \\
    \multicolumn{2}{l}{\universitylocation}
  \end{tabular}
}

\input{template}

\begin{document}
\templatePortrait
\templatePagecfg
\onehalfspacing

\begin{abstractd}
Este portafolio reúne las evidencias de la unidad de aprendizaje Álgebra Lineal en el programa de Ingeniero Agrónomo de la UANL. Cada actividad documenta su consigna, sustento teórico, producto, análisis y fuentes.
\end{abstractd}

\templateIndex
\templateFinalcfg

\section{Evidencias}
\begin{itemize}
  \item Actividad 1, Fase 1: mapa cognitivo de telaraña sobre aplicaciones del álgebra lineal en la vida cotidiana.
\end{itemize}

\bibliography{algebra-lineal}
\end{document}
